% Erratum E8 for inclusion in main.tex.
% Suggested use: % Erratum E8 for inclusion in main.tex.
% Suggested use: % Erratum E8 for inclusion in main.tex.
% Suggested use: % Erratum E8 for inclusion in main.tex.
% Suggested use: \input{erratum_E8_prop24_lower_bound}

\subsection*{Erratum E8: the lower-bound proof in Proposition 2.4}
\label{s:erratum:E8}

The lower-bound part of the proof of Proposition 2.4 invokes the
identity (9).  Since the $E$-part of (9) is false, that
part of the proof should either use only the corrected distance-function
identity from Erratum E3, or be replaced by a direct Hilbert-space duality
argument.  We record such a replacement here.

The first half of the proof in the paper already proves the inequality
``$\leq$'' in equation (36).  It remains only to justify the reverse
inequality in the interior case $0<\nu<\gamma$.  Put
\[
        \theta:=\frac{\nu}{\gamma},
        \qquad
        B:=(T^*T)^\gamma .
\]
For nonzero $\omega\in X$ define
\[
        F(x)
        :=
        \sup_{\omega\neq0}
        \frac{|\scalar{x,\omega}{}|}
             {\norm{\omega}{}^{1-\theta}\norm{B\omega}{}^\theta} .
\]
By homogeneity this is the same quantity as the left-hand side of
equation (36).  If $F(x)=\infty$, there is nothing to prove, so assume
$F(x)<\infty$.

We use the following dual representation of the $K$-functional.  From
equation (24), or equivalently from the spectral theorem,
\[
        |K_t^\gamma(x)|^2
        =
        \scalar{M_t x,x}{},
        \qquad
        M_t:=t^2(t^2+B^2)^{-1}.
\]
Since $B$ is bounded and self-adjoint, $M_t$ is bounded, positive, and
boundedly invertible, with
\[
        M_t^{-1}=I+t^{-2}B^2.
\]
Therefore
\begin{align}
        K_t^\gamma(x)
        &=
        \sup_{\omega\neq0}
        \frac{|\scalar{x,\omega}{}|}
             {\bigl(\scalar{M_t^{-1}\omega,\omega}{}\bigr)^{1/2}}        
        \notag\\
        &=
        \sup_{\omega\neq0}
        \frac{|\scalar{x,\omega}{}|}
             {\left(\norm{\omega}{}^2+t^{-2}\norm{B\omega}{}^2\right)^{1/2}} .
        \label{e:erratum:E8:dual-K}
\end{align}
Indeed, the identity is the elementary Hilbert-space formula
\[
        \scalar{M_t x,x}{}^{1/2}
        =
        \sup_{\omega\neq0}
        \frac{|\scalar{x,\omega}{}|}
             {\scalar{M_t^{-1}\omega,\omega}{}^{1/2}} .
\]

By the definition of $F(x)$,
\[
        |\scalar{x,\omega}{}|
        \leq
        F(x)\norm{\omega}{}^{1-\theta}\norm{B\omega}{}^\theta .
\]
Using Young's inequality in the form of equation (5), with
$a=\norm{\omega}{}^2$ and $b=t^{-2}\norm{B\omega}{}^2$, gives
\[
        N_\theta
        \norm{\omega}{}^{1-\theta}
        \bigl(t^{-1}\norm{B\omega}{}\bigr)^\theta
        \leq
        \left(\norm{\omega}{}^2+t^{-2}\norm{B\omega}{}^2\right)^{1/2}.
\]
Equivalently,
\[
        \norm{\omega}{}^{1-\theta}\norm{B\omega}{}^\theta
        \leq
        N_\theta^{-1} t^\theta
        \left(\norm{\omega}{}^2+t^{-2}\norm{B\omega}{}^2\right)^{1/2}.
\]
Substituting this into \eqref{e:erratum:E8:dual-K} yields, for every
$t>0$,
\[
        K_t^\gamma(x)
        \leq
        N_\theta^{-1} F(x) t^\theta .
\]
Consequently,
\[
        \norm{x}{\nu:\gamma}
        =
        \sup_{t>0}t^{-\theta}K_t^\gamma(x)
        \leq
        N_\theta^{-1}F(x).
\]
Thus
\[
        N_{\nu/\gamma}\norm{x}{\nu:\gamma}
        \leq
        F(x),
\]
which is precisely the missing inequality ``$\geq$'' in equation (36).
Combined with the already proved opposite inequality, this gives
\[
        F(x)=N_{\nu/\gamma}\norm{x}{\nu:\gamma}.
\]

This argument replaces the appeal to the variational-source-condition results
and avoids any use of the false $E$-identity in (9).  The endpoint
cases $\nu=0$ and $\nu=\gamma$ remain as stated in the original proof, subject
only to the restriction $\gamma>0$ noted in Erratum E7.


\subsection*{Erratum E8: the lower-bound proof in Proposition 2.4}
\label{s:erratum:E8}

The lower-bound part of the proof of Proposition 2.4 invokes the
identity (9).  Since the $E$-part of (9) is false, that
part of the proof should either use only the corrected distance-function
identity from Erratum E3, or be replaced by a direct Hilbert-space duality
argument.  We record such a replacement here.

The first half of the proof in the paper already proves the inequality
``$\leq$'' in equation (36).  It remains only to justify the reverse
inequality in the interior case $0<\nu<\gamma$.  Put
\[
        \theta:=\frac{\nu}{\gamma},
        \qquad
        B:=(T^*T)^\gamma .
\]
For nonzero $\omega\in X$ define
\[
        F(x)
        :=
        \sup_{\omega\neq0}
        \frac{|\scalar{x,\omega}{}|}
             {\norm{\omega}{}^{1-\theta}\norm{B\omega}{}^\theta} .
\]
By homogeneity this is the same quantity as the left-hand side of
equation (36).  If $F(x)=\infty$, there is nothing to prove, so assume
$F(x)<\infty$.

We use the following dual representation of the $K$-functional.  From
equation (24), or equivalently from the spectral theorem,
\[
        |K_t^\gamma(x)|^2
        =
        \scalar{M_t x,x}{},
        \qquad
        M_t:=t^2(t^2+B^2)^{-1}.
\]
Since $B$ is bounded and self-adjoint, $M_t$ is bounded, positive, and
boundedly invertible, with
\[
        M_t^{-1}=I+t^{-2}B^2.
\]
Therefore
\begin{align}
        K_t^\gamma(x)
        &=
        \sup_{\omega\neq0}
        \frac{|\scalar{x,\omega}{}|}
             {\bigl(\scalar{M_t^{-1}\omega,\omega}{}\bigr)^{1/2}}        
        \notag\\
        &=
        \sup_{\omega\neq0}
        \frac{|\scalar{x,\omega}{}|}
             {\left(\norm{\omega}{}^2+t^{-2}\norm{B\omega}{}^2\right)^{1/2}} .
        \label{e:erratum:E8:dual-K}
\end{align}
Indeed, the identity is the elementary Hilbert-space formula
\[
        \scalar{M_t x,x}{}^{1/2}
        =
        \sup_{\omega\neq0}
        \frac{|\scalar{x,\omega}{}|}
             {\scalar{M_t^{-1}\omega,\omega}{}^{1/2}} .
\]

By the definition of $F(x)$,
\[
        |\scalar{x,\omega}{}|
        \leq
        F(x)\norm{\omega}{}^{1-\theta}\norm{B\omega}{}^\theta .
\]
Using Young's inequality in the form of equation (5), with
$a=\norm{\omega}{}^2$ and $b=t^{-2}\norm{B\omega}{}^2$, gives
\[
        N_\theta
        \norm{\omega}{}^{1-\theta}
        \bigl(t^{-1}\norm{B\omega}{}\bigr)^\theta
        \leq
        \left(\norm{\omega}{}^2+t^{-2}\norm{B\omega}{}^2\right)^{1/2}.
\]
Equivalently,
\[
        \norm{\omega}{}^{1-\theta}\norm{B\omega}{}^\theta
        \leq
        N_\theta^{-1} t^\theta
        \left(\norm{\omega}{}^2+t^{-2}\norm{B\omega}{}^2\right)^{1/2}.
\]
Substituting this into \eqref{e:erratum:E8:dual-K} yields, for every
$t>0$,
\[
        K_t^\gamma(x)
        \leq
        N_\theta^{-1} F(x) t^\theta .
\]
Consequently,
\[
        \norm{x}{\nu:\gamma}
        =
        \sup_{t>0}t^{-\theta}K_t^\gamma(x)
        \leq
        N_\theta^{-1}F(x).
\]
Thus
\[
        N_{\nu/\gamma}\norm{x}{\nu:\gamma}
        \leq
        F(x),
\]
which is precisely the missing inequality ``$\geq$'' in equation (36).
Combined with the already proved opposite inequality, this gives
\[
        F(x)=N_{\nu/\gamma}\norm{x}{\nu:\gamma}.
\]

This argument replaces the appeal to the variational-source-condition results
and avoids any use of the false $E$-identity in (9).  The endpoint
cases $\nu=0$ and $\nu=\gamma$ remain as stated in the original proof, subject
only to the restriction $\gamma>0$ noted in Erratum E7.


\subsection*{Erratum E8: the lower-bound proof in Proposition 2.4}
\label{s:erratum:E8}

The lower-bound part of the proof of Proposition 2.4 invokes the
identity (9).  Since the $E$-part of (9) is false, that
part of the proof should either use only the corrected distance-function
identity from Erratum E3, or be replaced by a direct Hilbert-space duality
argument.  We record such a replacement here.

The first half of the proof in the paper already proves the inequality
``$\leq$'' in equation (36).  It remains only to justify the reverse
inequality in the interior case $0<\nu<\gamma$.  Put
\[
        \theta:=\frac{\nu}{\gamma},
        \qquad
        B:=(T^*T)^\gamma .
\]
For nonzero $\omega\in X$ define
\[
        F(x)
        :=
        \sup_{\omega\neq0}
        \frac{|\scalar{x,\omega}{}|}
             {\norm{\omega}{}^{1-\theta}\norm{B\omega}{}^\theta} .
\]
By homogeneity this is the same quantity as the left-hand side of
equation (36).  If $F(x)=\infty$, there is nothing to prove, so assume
$F(x)<\infty$.

We use the following dual representation of the $K$-functional.  From
equation (24), or equivalently from the spectral theorem,
\[
        |K_t^\gamma(x)|^2
        =
        \scalar{M_t x,x}{},
        \qquad
        M_t:=t^2(t^2+B^2)^{-1}.
\]
Since $B$ is bounded and self-adjoint, $M_t$ is bounded, positive, and
boundedly invertible, with
\[
        M_t^{-1}=I+t^{-2}B^2.
\]
Therefore
\begin{align}
        K_t^\gamma(x)
        &=
        \sup_{\omega\neq0}
        \frac{|\scalar{x,\omega}{}|}
             {\bigl(\scalar{M_t^{-1}\omega,\omega}{}\bigr)^{1/2}}        
        \notag\\
        &=
        \sup_{\omega\neq0}
        \frac{|\scalar{x,\omega}{}|}
             {\left(\norm{\omega}{}^2+t^{-2}\norm{B\omega}{}^2\right)^{1/2}} .
        \label{e:erratum:E8:dual-K}
\end{align}
Indeed, the identity is the elementary Hilbert-space formula
\[
        \scalar{M_t x,x}{}^{1/2}
        =
        \sup_{\omega\neq0}
        \frac{|\scalar{x,\omega}{}|}
             {\scalar{M_t^{-1}\omega,\omega}{}^{1/2}} .
\]

By the definition of $F(x)$,
\[
        |\scalar{x,\omega}{}|
        \leq
        F(x)\norm{\omega}{}^{1-\theta}\norm{B\omega}{}^\theta .
\]
Using Young's inequality in the form of equation (5), with
$a=\norm{\omega}{}^2$ and $b=t^{-2}\norm{B\omega}{}^2$, gives
\[
        N_\theta
        \norm{\omega}{}^{1-\theta}
        \bigl(t^{-1}\norm{B\omega}{}\bigr)^\theta
        \leq
        \left(\norm{\omega}{}^2+t^{-2}\norm{B\omega}{}^2\right)^{1/2}.
\]
Equivalently,
\[
        \norm{\omega}{}^{1-\theta}\norm{B\omega}{}^\theta
        \leq
        N_\theta^{-1} t^\theta
        \left(\norm{\omega}{}^2+t^{-2}\norm{B\omega}{}^2\right)^{1/2}.
\]
Substituting this into \eqref{e:erratum:E8:dual-K} yields, for every
$t>0$,
\[
        K_t^\gamma(x)
        \leq
        N_\theta^{-1} F(x) t^\theta .
\]
Consequently,
\[
        \norm{x}{\nu:\gamma}
        =
        \sup_{t>0}t^{-\theta}K_t^\gamma(x)
        \leq
        N_\theta^{-1}F(x).
\]
Thus
\[
        N_{\nu/\gamma}\norm{x}{\nu:\gamma}
        \leq
        F(x),
\]
which is precisely the missing inequality ``$\geq$'' in equation (36).
Combined with the already proved opposite inequality, this gives
\[
        F(x)=N_{\nu/\gamma}\norm{x}{\nu:\gamma}.
\]

This argument replaces the appeal to the variational-source-condition results
and avoids any use of the false $E$-identity in (9).  The endpoint
cases $\nu=0$ and $\nu=\gamma$ remain as stated in the original proof, subject
only to the restriction $\gamma>0$ noted in Erratum E7.


\subsection*{Erratum E8: the lower-bound proof in Proposition 2.4}
\label{s:erratum:E8}

The lower-bound part of the proof of Proposition 2.4 invokes the
identity (9).  Since the $E$-part of (9) is false, that
part of the proof should either use only the corrected distance-function
identity from Erratum E3, or be replaced by a direct Hilbert-space duality
argument.  We record such a replacement here.

The first half of the proof in the paper already proves the inequality
``$\leq$'' in equation (36).  It remains only to justify the reverse
inequality in the interior case $0<\nu<\gamma$.  Put
\[
        \theta:=\frac{\nu}{\gamma},
        \qquad
        B:=(T^*T)^\gamma .
\]
For nonzero $\omega\in X$ define
\[
        F(x)
        :=
        \sup_{\omega\neq0}
        \frac{|\scalar{x,\omega}{}|}
             {\norm{\omega}{}^{1-\theta}\norm{B\omega}{}^\theta} .
\]
By homogeneity this is the same quantity as the left-hand side of
equation (36).  If $F(x)=\infty$, there is nothing to prove, so assume
$F(x)<\infty$.

We use the following dual representation of the $K$-functional.  From
equation (24), or equivalently from the spectral theorem,
\[
        |K_t^\gamma(x)|^2
        =
        \scalar{M_t x,x}{},
        \qquad
        M_t:=t^2(t^2+B^2)^{-1}.
\]
Since $B$ is bounded and self-adjoint, $M_t$ is bounded, positive, and
boundedly invertible, with
\[
        M_t^{-1}=I+t^{-2}B^2.
\]
Therefore
\begin{align}
        K_t^\gamma(x)
        &=
        \sup_{\omega\neq0}
        \frac{|\scalar{x,\omega}{}|}
             {\bigl(\scalar{M_t^{-1}\omega,\omega}{}\bigr)^{1/2}}        
        \notag\\
        &=
        \sup_{\omega\neq0}
        \frac{|\scalar{x,\omega}{}|}
             {\left(\norm{\omega}{}^2+t^{-2}\norm{B\omega}{}^2\right)^{1/2}} .
        \label{e:erratum:E8:dual-K}
\end{align}
Indeed, the identity is the elementary Hilbert-space formula
\[
        \scalar{M_t x,x}{}^{1/2}
        =
        \sup_{\omega\neq0}
        \frac{|\scalar{x,\omega}{}|}
             {\scalar{M_t^{-1}\omega,\omega}{}^{1/2}} .
\]

By the definition of $F(x)$,
\[
        |\scalar{x,\omega}{}|
        \leq
        F(x)\norm{\omega}{}^{1-\theta}\norm{B\omega}{}^\theta .
\]
Using Young's inequality in the form of equation (5), with
$a=\norm{\omega}{}^2$ and $b=t^{-2}\norm{B\omega}{}^2$, gives
\[
        N_\theta
        \norm{\omega}{}^{1-\theta}
        \bigl(t^{-1}\norm{B\omega}{}\bigr)^\theta
        \leq
        \left(\norm{\omega}{}^2+t^{-2}\norm{B\omega}{}^2\right)^{1/2}.
\]
Equivalently,
\[
        \norm{\omega}{}^{1-\theta}\norm{B\omega}{}^\theta
        \leq
        N_\theta^{-1} t^\theta
        \left(\norm{\omega}{}^2+t^{-2}\norm{B\omega}{}^2\right)^{1/2}.
\]
Substituting this into \eqref{e:erratum:E8:dual-K} yields, for every
$t>0$,
\[
        K_t^\gamma(x)
        \leq
        N_\theta^{-1} F(x) t^\theta .
\]
Consequently,
\[
        \norm{x}{\nu:\gamma}
        =
        \sup_{t>0}t^{-\theta}K_t^\gamma(x)
        \leq
        N_\theta^{-1}F(x).
\]
Thus
\[
        N_{\nu/\gamma}\norm{x}{\nu:\gamma}
        \leq
        F(x),
\]
which is precisely the missing inequality ``$\geq$'' in equation (36).
Combined with the already proved opposite inequality, this gives
\[
        F(x)=N_{\nu/\gamma}\norm{x}{\nu:\gamma}.
\]

This argument replaces the appeal to the variational-source-condition results
and avoids any use of the false $E$-identity in (9).  The endpoint
cases $\nu=0$ and $\nu=\gamma$ remain as stated in the original proof, subject
only to the restriction $\gamma>0$ noted in Erratum E7.
